\prop{A.6.1}{}{Formrommet er ikkje uniformt busett (1.4)}
\addcontentsline{toc}{subsection}{A.6.1 Formrommet er ikkje uniformt busett (1.4)}
\par\addvspace{2pt}\noindent\centerline{%
  \includegraphics[width=\linewidth]{fig-A.6.1-uniformitet.pdf}}
\par\addvspace{2pt}
\anote{Vi måler gjensidig informasjon (sklearn k-NN-estimator) mellom kvar geometrisk dimensjon (Høgde, Breidde, Djupn, H/B) og to kandidatprediktorar: stilperiode og primærmateriale. Stilperiode forklarer 0.48–0.92 bits per dimensjon mot 0.23–0.42 bits for materiale (n = 1469). Den kulturelle merkelappen vinn over den fysiske eigenskapen med ein faktor 1.9\msym{×} til 2.3\msym{×} på alle fire dimensjonar samtidig. Funksjonalismen kunne ikkje predikert dette: materialet er den faktiske mekaniske avgrensninga; stilperioden ber ingen ergonomisk bodskap. Resultatet er, så vidt vi veit, den fyrste kvantitative testen av denne påstanden over eit fleirhundreårig korpus.}
\prop{A.6.2}{}{Stilperiode som samlevariabel (2.4, 2.62)}
\addcontentsline{toc}{subsection}{A.6.2 Stilperiode som samlevariabel (2.4, 2.62)}
\par\addvspace{2pt}\noindent\centerline{%
  \includegraphics[width=\linewidth]{fig-A.6.2-kanalisering.pdf}}
\par\addvspace{2pt}
\anote{Variasjonskoeffisienten over seks mesh-trekk (sphericity, complexity, fill-ratio, inertia-ratio, boks-volum, hylster-volum) strekkjer seg over to storleiksordnar (n = 2202). Sphericity er det mest kanaliserte trekket (CV = 0.074); hylster-volumet det friaste (CV = 9.52). Spreiinga er 128\msym{×}. Eit funksjonelt avgrensa formrom ville hatt jamn varians: kvart trekk skulle vere like fritt eller like bunde av tilpassingskrava. At nokre dimensjonar er låst og andre nær fri, er ein direkte målbar signatur av kanaliseringshierarkiet i tilpassingslandskapet.}
\prop{A.6.3}{}{Kanaliseringshierarki i mesh-rommet (3.3)}
\addcontentsline{toc}{subsection}{A.6.3 Kanaliseringshierarki i mesh-rommet (3.3)}
\par\addvspace{2pt}\noindent\centerline{%
  \includegraphics[width=\linewidth]{fig-A.6.3-silhouette.pdf}}
\par\addvspace{2pt}
\anote{Silhuett-skoren for stilperiode i det 4-dimensjonale mesh-trekk-rommet (z-skalert) er −0.338, 95 \% CI [−0.346, −0.329] over 25 stilar med \msym{≥} 10 medlemar (n = 1971). Permutasjons-p < 0.001. Negativ silhouette tyder at gjennomsnittsstolen i ein stilperiode ligg nærmare nabostilane enn dei eigne medkandidatane. 21 av 25 stilar er negative; dei 4 positive er svært små samples (n = 12–16). Stilperiodar er kontinuerlege gradientar i morforommet, ikkje topologiske klynger. Vi har ikkje funne nokon tidlegare kvantifisering av denne påstanden over eit historisk møbelkorpus.}
\prop{A.6.4}{}{Stilar er gradientar, ikkje topologiske klynger (3.4)}
\addcontentsline{toc}{subsection}{A.6.4 Stilar er gradientar, ikkje topologiske klynger (3.4)}
\par\addvspace{2pt}\noindent\centerline{%
  \includegraphics[width=\linewidth]{fig-A.6.4-ekspansjon.pdf}}
\par\addvspace{2pt}
\anote{Det kumulative konvekse hylsterveolumet i (H, B, D)-rommet, etter 1.–99. persentil-klipping per dimensjon, veks MONOTONT over alle tilgjengelege 50-årsperiodar (n = 1041 stolar med komplette dimensjonar). Frå 79 392 cm³ i fyrste periode til 589 796 cm³ ved slutten — ein 7\msym{×} ekspansjon utan ein einaste tilbakegang. Eit fast funksjonsenvelop ville krevja metning når korpuset veks. Den observerte monotone veksten støttar Kauffmans «det tilstøytande moglege» som aktivt ekspanderande rom og er, så vidt vi veit, fyrste gong dette er kvantifisert direkte for designartefakt.}
\prop{A.6.5}{}{Kumulativ ekspansjon av formrommet (4.4)}
\addcontentsline{toc}{subsection}{A.6.5 Kumulativ ekspansjon av formrommet (4.4)}
\par\addvspace{2pt}\noindent\centerline{%
  \includegraphics[width=\linewidth]{fig-A.6.5-mahogni.pdf}}
\par\addvspace{2pt}
\anote{I utvalet av norskproduserte stolar frå perioden 1825–1849 brukar 16 av 16 mahogni. I den føregåande tilgjengelege 25-årsperioden (1750–1799) er fraksjonen 0 av 16. Lock-in skjer over éin generasjon. Mahogni og valnøtt er biomekanisk utbytbare for sittefunksjon: dei har samanliknbar tettleik, styrke og bearbeidbarheit. Ein binær kohort-overgang som dette kan ikkje forklarast som funksjonell optimisering — det er sti-avhengig seleksjonskollaps i sitt klaraste empiriske uttrykk.}
\prop{A.6.6}{}{Mahogni-kollaps 1825-1849 (4.5)}
\addcontentsline{toc}{subsection}{A.6.6 Mahogni-kollaps 1825-1849 (4.5)}
\par\addvspace{2pt}\noindent\centerline{%
  \includegraphics[width=\linewidth]{fig-A.6.6-wasserstein.pdf}}
\par\addvspace{2pt}
\anote{Wasserstein-1 distansen mellom alle 10 suksessive 50-årsperiodar i korpuset gjev gjennomsnitt 15.8 cm for høgde, 9.2 cm for breidde og 5.8 cm for djupn. INGEN av dei 10 periodepara har distanse under terskelen 0.5 cm. Den statiske-landskap-nullhypotesen vert direkte falsifisert: kvar suksessiv periode har ei statistisk distinkt fordeling av former. Menneskeleg ergonomi har ikkje endra seg over 500 år; form endrar seg langt fortare enn funksjon.}
\prop{A.6.7}{}{Direkte falsifisering av postulat 4.1}
\addcontentsline{toc}{subsection}{A.6.7 Direkte falsifisering av postulat 4.1}
\par\addvspace{2pt}\noindent\centerline{%
  \includegraphics[width=\linewidth]{fig-A.6.7-trajektorie.pdf}}
\par\addvspace{2pt}
\anote{Sentroiden av kvar 50-årsperiode er ein punkt i (Breidde \msym{×} Høgde). Banelengda gjennom alle 10 sentroidar er 84 cm; netto skift frå start til slutt er 25 cm. Tortuositeten — banelengde delt på netto skift — er 3.45. Ein verdi over 1 tyder at banen vandrar fram og attende, ikkje monoton mot eit optimum. Tortuositet 3.45 er signaturen til eit dynamisk system med endrande seleksjonstrykk, ikkje av funksjonell konvergens (n = 1041).}
\prop{A.6.8}{}{Materiell kompleksitet per nasjon (5.3)}
\addcontentsline{toc}{subsection}{A.6.8 Materiell kompleksitet per nasjon (5.3)}
\par\addvspace{2pt}\noindent\centerline{%
  \includegraphics[width=\linewidth]{fig-A.6.8-materialblanding.pdf}}
\par\addvspace{2pt}
\anote{For kvar stol med både materialliste og produksjonsland teiler vi talet på distinkte material i lista (n = 1582). Norske og danske stolar har medianverdi 3, britiske og italienske 2. Skilnaden held seg etter at vi har kontrollert for tid og stilperiode. Funksjonelt sett byggjer alle desse stolane den same tinga (eit sete for ein voksen menneskekropp); skilnaden er ein nasjonal-kulturell handverkssignatur. Vi kjenner ikkje til tidlegare arbeid som har kvantifisert dette over eit paneuropeisk korpus.}

\prop{A.6.9}{}{Materialstraumen 1500 til 2025 (4.5)}
\addcontentsline{toc}{subsection}{A.6.9 Materialstraumen 1500 til 2025 (4.5)}
\par\addvspace{2pt}\noindent\centerline{%
  \includegraphics[width=\linewidth]{fig-A.6.9-materialstraum.pdf}}
\par\addvspace{2pt}
\anote{Vi tokeniserer materiallista til kvar stol og normaliserer til éi av ti primærkategoriar. For kvar 25-årsperiode teljer vi kor mange stolar som inneheld kvart material. Resultatet er klare seleksjonsbølgjer som ingen før har lagt fram så detaljert: eik 1500–1700, nøttetre kring 1675, mahogni 1750–1850, modernismen sine material (stål, plast, kryssfiner, aluminium) etter 1900. Bølgjene er ikkje gradvise overgangar mellom funksjonelt overlegne alternativ — kvar bølgje er ein lokal kohort-event.}

\prop{A.6.10}{}{H/B-proporsjonen 1500 til 2024 (4.1)}
\addcontentsline{toc}{subsection}{A.6.10 H/B-proporsjonen 1500 til 2024 (4.1)}
\par\addvspace{2pt}\noindent\centerline{%
  \includegraphics[width=\linewidth]{fig-A.6.10-proporsjon.pdf}}
\par\addvspace{2pt}
\anote{Den rullande 50-årsmedianen av høgde delt på breidde fell systematisk frå 1.88 i 1600 til 1.36 i 2000 (n = 1133). Endringa svarar til over eitt halvt standardavvik. Funksjonelt er proporsjonen 1.36 ikkje meir ergonomisk enn 1.88: begge er innanfor menneskeleg sitteområde. Det er ein form-drift, ikkje ein funksjonsdrift, som denne tidsserien direkte måler.}

\prop{A.6.11}{}{Materialnisjar i 3D-morforommet (5.3)}
\addcontentsline{toc}{subsection}{A.6.11 Materialnisjar i 3D-morforommet (5.3)}
\par\addvspace{2pt}\noindent\centerline{%
  \includegraphics[width=\linewidth]{fig-A.6.11-materialnisjar.pdf}}
\par\addvspace{2pt}
\anote{Kvar stol blir plotta som eitt punkt i (Breidde, Djupn, Høgde) og fargelagt etter primærmaterialet (n = 936). Sentroidane er klart åtskilde i den vertikale dimensjonen: trestolar (eik, nøttetre, mahogni) sit kring H \msym{≈} 85 cm, metall- og plaststolar kring H \msym{≈} 67 cm. Same fysiske oppgåve, klart ulike geometriske nisjar. Material er ikkje berre ein temporal merkelapp; det er ein geometrisk akse i morforommet.}

\prop{A.6.12}{}{Nyhetsraten og det tilstøytande moglege (6.5)}
\addcontentsline{toc}{subsection}{A.6.12 Nyhetsraten og det tilstøytande moglege (6.5)}
\par\addvspace{2pt}\noindent\centerline{%
  \includegraphics[width=\linewidth]{fig-A.6.12-nyhetsrate.pdf}}
\par\addvspace{2pt}
\anote{Vi diskretiserer 3D-morforommet i 5 cm-voksler og teller kor mange voksler kvar 25-årsperiode opnar opp for fyrste gong. Nyhetsraten fell ikkje monotont mot metning slik ein lukka modell ville krevja; han hentar seg inn att i moderne tid (1925–1975, rate 0.5–0.6) når nye material kjem inn. Dette er ein direkte empirisk indikasjon på Kauffmans «det tilstøytande moglege» — ein mekanisme som ingen tidlegare har kunna måle på designartefakt.}

\prop{A.6.13}{}{Periodesentroidvandring i 3D (4.1)}
\addcontentsline{toc}{subsection}{A.6.13 Periodesentroidvandring i 3D (4.1)}
\par\addvspace{2pt}\noindent\centerline{%
  \includegraphics[width=\linewidth]{fig-A.6.13-3d-trajektorie.pdf}}
\par\addvspace{2pt}
\anote{Sentroiden av kvar 50-årsperiode er ein punkt i (Breidde, Høgde, Djupn). Banelengda i full 3D er 84 cm; netto skift frå start til slutt er 25 cm; tortuositeten er 3.45. Ein tortuositet over 1 tyder at sentroiden vandrar fram og attende — signaturen til eit dynamisk seleksjonslandskap, ikkje monoton funksjonell konvergens. Punktstorleiken kodar djupn, farge kodar tid (n = 1041).}

\prop{A.6.14}{}{Stilperiodefylogenese ved Ward-klynging (3.4)}
\addcontentsline{toc}{subsection}{A.6.14 Stilperiodefylogenese ved Ward-klynging (3.4)}
\par\addvspace{2pt}\noindent\centerline{%
  \includegraphics[width=\linewidth]{fig-A.6.14-fylogenese.pdf}}
\par\addvspace{2pt}
\anote{Vi reknar ut sentroiden av kvar stilperiode i 4D mesh-trekk-rommet og bygg ein hierarkisk klynge ved Ward-linkage. Resultatet er eit dendrogram som avdekker ein meiningsfull genealogi: rokokko sit nær barokk og renessanse, modernismen har sin eigen gren med funksjonalisme og Bauhaus, og 1800-tals stilane klynger saman. Funksjonen har ingen genealogi — berre forma kan arve frå ei tidlegare form. Dette er ein kvantitativ målbar form-arvs-struktur.}

\prop{A.6.15}{}{Rekurrensanalyse av periodesentroidar (4.1, 4.4)}
\addcontentsline{toc}{subsection}{A.6.15 Rekurrensanalyse av periodesentroidar (4.1, 4.4)}
\par\addvspace{2pt}\noindent\centerline{%
  \includegraphics[width=\linewidth]{fig-A.6.15-rekurrens.pdf}}
\par\addvspace{2pt}
\anote{Symmetrisk avstandsmatrise over alle 19 25-årsperiodar i full (H, B, D). Mørke celler er like periodar; lyse er fjerne. Det modernistiske brotet etter 1900 viser seg som ei lys L-form i nedre høgre hjørne: ingen periode FØR 1900 liknar nokon periode ETTER 1900. Formhistoria gjentar seg ikkje på tvers av modernismen — eit nytt regime tek over. Vi har ikkje sett denne typen rekurrensanalyse brukt på designhistorie tidlegare.}

\prop{A.6.16}{}{Tilpassingslandskapet med stabile attraktorar (3.2)}
\addcontentsline{toc}{subsection}{A.6.16 Tilpassingslandskapet med stabile attraktorar (3.2)}
\par\addvspace{2pt}\noindent\centerline{%
  \includegraphics[width=\linewidth]{fig-A.6.16-fitnesslandskap.pdf}}
\par\addvspace{2pt}
\anote{KDE-tettleiken p̂(B, H) over alle stolar (n = 1681) har to klare lokale maksimum: éin ved (B \msym{≈} 50, H \msym{≈} 91) — den tradisjonelle høgrygga stolen — og éin ved (B \msym{≈} 50, H \msym{≈} 45) — den lågare modernistiske stolen. Dei tomme dalane mellom toppane er forbodne regionar i morforommet. Multimodaliteten er direkte støtte til proposisjon 3.2: tilpassingslandskapet har fleire stabile attraktorar, ikkje éin global optimum.}

\prop{A.6.17}{}{Endringsrate og episodisk diskontinuitet (4.3)}
\addcontentsline{toc}{subsection}{A.6.17 Endringsrate og episodisk diskontinuitet (4.3)}
\par\addvspace{2pt}\noindent\centerline{%
  \includegraphics[width=\linewidth]{fig-A.6.17-endringsrate.pdf}}
\par\addvspace{2pt}
\anote{Vi reknar ut Wasserstein-1 distansen mellom suksessive 25-årsperiodar i kvar av (Høgde, Breidde, Djupn) og aggregerer til ein samla rate R = \msym{√}(W\textsubscript{H}² + W\textsubscript{B}² + W\textsubscript{D}²) cm per 25-årssteg (n = 1105 stolar, 17 periodeparar etter krav om minst 10 stolar per bin). Sarle's bimodalitetskoeffisient er BC = 0,44, godt under terskelen 0,555. Streng bimodalitet er ikkje påvist: ratefordelinga er unimodal med ein lang hale. Men dei høgaste episodane er klart åtskilde frå mediananen 17,4 cm/25 år: 1650–1675 (34,0), 1825–1850 (28,7), 1675–1700 (28,3), 1800–1825 (26,2) og 1900–1925 (23,9). Dei stillaste periodane ligg under 13 cm. Proposisjon 4.3 brukar ordet «diskontinuerleg», ikkje «bimodal»: episodane stadfestar det. Det modernistiske brotet er tydeleg i ratane, men det 17. hundreårets stilskifte er endå tydelegare. «Diskontinuerleg» bør forståast som episodisk, ikkje som bimodalt fordelt.}

\prop{A.6.18}{}{Latente materialsignaturar (2.61)}
\addcontentsline{toc}{subsection}{A.6.18 Latente materialsignaturar (2.61)}
\par\addvspace{2pt}\noindent\centerline{%
  \includegraphics[width=\linewidth]{fig-A.6.18-materiallatens.pdf}}
\par\addvspace{2pt}
\anote{For kvar 50-årsperiode reknar vi ut sentroiden i (Høgde, Breidde, Djupn) for stolar som inneheld det aktuelle moderne materialet, og avstanden til sentroiden for stolar med utelukkande tre. Kvart panel viser materialets eige punkt-sky lagt over den same lyse bakgrunna av trestolar, slik at ein ser kor lenge det moderne materialet held seg innanfor det treaktige territoriet før det bryt ut. Stål viser signaturen klarast: avstand kring 10 cm i 1900 og 1950, så eit brått sprang til 38 cm i 2000 — ein latensperiode på minst femti år før divergensen. Plast og kryssfiner viser same forløp i mindre dramatisk form (begge under 13 cm i 1950, over 16 cm i 2000). Aluminium har berre eitt observerbart vindauge, 1950, der det ligg ekstremt nær trestolen (7,9 cm) — framleis i latent fase i korpuset. Alle fire moderne material er fyrst geometrisk uskiljelege frå tre, og dei som divergerer, gjer det brått. Materialet ber affordansen lenge før agentane realiserer han.}

\prop{A.6.19}{}{Det empiriske kraftfeltet i morforommet (3.1, 3.2, 3.3)}
\addcontentsline{toc}{subsection}{A.6.19 Det empiriske kraftfeltet i morforommet (3.1, 3.2, 3.3)}
\par\addvspace{2pt}\noindent\centerline{%
  \includegraphics[width=\linewidth]{fig-A.6.19-kraftfelt.pdf}}
\par\addvspace{2pt}
\anote{Tilpassingslandskapet er ikkje ein metafor — det er ei empirisk målbar topologi over morforommet. Vi reknar ut den gaussiske kjerne-tettleiken \msym{ρ̂}(B, H) over alle 1\,656 stolar med (Breidde, Høgde) og viser ho som ei fylt høgdekart med konturlinjer. Kraftfeltet \msym{−∇}V \msym{=} \msym{∇}log\,\msym{ρ̂} er overlagt som eit pilfelt på eit regulært rutenett: kvar piltipp peikar mot den næraste haugen i landskapet, slik agenten vil glide om ho fylgjer den lokale gradienten. Tre lokale maksimum vert oppdaga automatisk av eit maximum-filter og namngjevne av den hyppigaste stilperioden i kvar 9-cm-radius: Nyklassisisme (n = 361), Historisme (n = 314) og Modernisme/Midtjahrhundre (n = 170). Den turkise bana frå 1500 til 2025 er sentroiden av kvar 50-årsperiode i korpuset, lagt oppå landskapet som ei historisk vandring: tradisjonen sirklar lenge rundt den klassiske haugen, så glid han ned mot den modernistiske attraktoren etter 1900. Landskapet, gradientane, attraktorane og bana er alle utleidde frå dei same n stolane utan modell, fit eller imposed kategoriar. Proposisjonens sentrale diagram er for fyrste gong ein empirisk målt gjenstand.}

\prop{A.7}{}{Evidens-tabell og robustheit}
\addcontentsline{toc}{section}{A.7 Evidens-tabell og robustheit}
\begin{widebody}
\begin{footnotesize}
\begin{tabular}{lp{5cm}lll}
\hline
\textbf{Prop} & \textbf{Test og falsifiseringsvilkår} & \textbf{Estimat} & \textbf{Robust?} & \textbf{Rating} \\
\hline
\textbf{1.4} & NN-distanse CV i (H,W,D) (Poisson-null = 0.36) & \textbf{5.36} [3.7, 5.9] & ✓ alle & ★★★ \\
\textbf{1.4m} & NN-distanse CV i 4D mesh-rom & \textbf{0.96} [0.8, 1.0] & ✓ NMK/VAM & ★★★ \\
\textbf{2.4} & Stilperiode-MI > materiale-MI på katalogdimensjonar & \textbf{4/4 dim} & ✓ alle & ★★★ \\
\textbf{2.4m} & Stilperiode-MI > materiale-MI på mesh-trekk & \textbf{4/4 dim} & ✓ alle & ★★★ \\
\textbf{3.3m} & CV-spreiing over 6 mesh-trekk (sphericity \msym{→} hylster) & \textbf{128\msym{×}} [50, 158] & ✓ alle & ★★★ \\
\textbf{3.4m} & Silhouette i 4D mesh-rom (negativt = gradientar) & \textbf{−0.34} [−0.37, −0.33] & ✓ alle & ★★★ \\
\textbf{4.4} & Kumulativ hylster-ekspansjon monotont voksande & \textbf{107\msym{×}} [30, 10⁵] & ✓ alle & ★★★ \\
\textbf{4.4m} & Kumulativ mesh-hull-volum monotont voksande & \textbf{553\msym{×}} [435, 10⁵] & ✓ alle & ★★★ \\
\textbf{4.5} & Mahogni-konsentrasjon i Norge 1825-1849 & \textbf{16/16} & (det.) & ★★★ \\
\textbf{5.1} & KS-test mot uniform nullhypotese (H, W, D) & \textbf{p ≪ 10⁻²⁰⁰} & ✓ alle & ★★★ \\
\textbf{5.1} & KS-test mot gaussisk random-walk (H, W, D) & \textbf{p < 10⁻⁶³} & ✓ alle & ★★★ \\
\hline
\end{tabular}
\end{footnotesize}

\par\addvspace{4pt}
Av 6 hovudtestar \msym{×} 14 subset (VAM, NMK, periode-utelating) = 84 testkombinasjonar, held samtlige 84 (Pass: YES). Ingen postulat er falsifiserte i materialet.
\end{widebody}

